\documentclass[10pt]{ijgtp}
\usepackage{cmap}
\usepackage[utf8]{inputenc}
\usepackage[T1]{fontenc}
\usepackage[english]{babel}

\newcommand{\cmd}[1]{\texttt{\detokenize{#1}}}

\begin{document}
\papertype{User Guide}
\title{IJGTP Class Documentation}
\begin{abstract}
This guide describes the public interface of \texttt{ijgtp.cls}: preparing
one manuscript or a collection of papers, setting author and publication
metadata, using figures and tables, and compiling bibliographies. It describes
the main commands available to authors, selected standard commands
whose behavior the class changes, and the additional table column types.
Internal commands containing \texttt{@} are implementation details and
are not intended for use in manuscripts.
\end{abstract}
\keywords{LaTeX; IJGTP; document class; author instructions}
\maketitle
\tableofcontents

\section{Files, packages, and class options}
\label{sec:setup}
Place \texttt{ijgtp.cls} and \texttt{scilight.bst} in the working directory
or in a directory searched by TeX. The title and footer require the image
files \texttt{header}, \texttt{F1}, and \texttt{cc-by}; the ORCID command
also requires \texttt{orcid}. Supply their PDF versions for PDFLaTeX and
their EPS versions for the LaTeX--dvips workflow. The sample plot files
are examples, not class dependencies.

The class is based on \texttt{article}, defaults to 10-point text, and
sets A4 journal geometry. Options are passed to \texttt{article}: for
example, \cmd{\documentclass[12pt]{ijgtp}} selects 12-point text.

Load optional packages explicitly. For example:
\begin{samepage}
\begin{verbatim}
\documentclass[10pt]{ijgtp}
\usepackage{amsmath,amssymb,cmap}
\usepackage[utf8]{inputenc}
\usepackage[T1]{fontenc}
\usepackage[english]{babel}
% \usepackage{float} % only if its facilities are needed
% \usepackage{epstopdf} \epstopdfsetup{suffix=} % convert eps figures in pdf
\end{verbatim}
\end{samepage}

The class already loads \texttt{natbib}, \texttt{bibunits},
\texttt{etoolbox}, \texttt{geometry}, \texttt{xcolor},
\texttt{microtype}, \texttt{setspace}, \texttt{times},
\texttt{titlesec}, \texttt{array}, \texttt{lineno},
\texttt{graphicx}, \texttt{caption}, \texttt{fancyhdr},
\texttt{lastpage}, \texttt{catchfile}, and \texttt{hyperref}.
Their commands remain available.

\section{A minimal manuscript}
\label{sec:minimal}
Define the metadata before \cmd{\maketitle}. Each \cmd{\author} starts
a new author record. The following affiliation, email, ORCID, and
correspondence commands apply to that author.
Avoid verbatim material inside the captured abstract environment.
\begin{verbatim}
\documentclass{ijgtp}
\begin{document}
\title{Title of the manuscript}
\author{Wednesday Addams}
\affiliation{Department of Physics, Nevermore Academy}
\email{inmemoryofwednesday@gmail.com}
\corresponding
\begin{abstract}
An overview of the question, method, and results.
\end{abstract}
\keywords{gravitation; theoretical physics}
\maketitle
\section{Introduction}\label{sec:intro}
The manuscript starts here.
\section{Conclusions}
See Section~\ref{sec:intro}.
\end{document}
\end{verbatim}
\cmd{template.tex} provides an extended example that can serve as a starting point for a manuscript.

\section{Author-facing command reference}
\label{sec:commands}
The following list covers the public commands and environments defined
or given paper-specific behavior by the class. Arguments in examples
are placeholders unless stated otherwise.

\subsection{Basic information}
\begin{description}
\item[\cmd{\title{title}}]
Sets the title for the next \cmd{\maketitle}. The title is retained
after printing, so it can be reused if no replacement is supplied.
Every manuscript is expected to have a title.
\item[\cmd{\papertype{type}}]
Sets the designation above the title, such as \texttt{Article} or
\texttt{Review}. The default is \texttt{Article}.
\item[\cmd{\author{full name}}]
Appends an author. Repeat this command for additional authors.
Author notes can be supplied with \cmd{\thanks{note}} inside the name.
\item[\cmd{\authorname{surname}{initials}}]
Overrides the current author's name parts used in running heads and
the suggested citation. It does not alter the displayed full name.
For example, use \cmd{\authorname{van Beethoven}{L.}} after
\cmd{\author{Ludwig van Beethoven}}. Without an override, the last
space-separated name component is treated as the surname and earlier
components supply initials. Use the override for compound surnames,
suffixes, or other names that this simple rule cannot represent.
\item[\cmd{\affiliation{address}}]
Adds an affiliation to the current author. Repeat it for multiple
affiliations. Identical affiliation text is shared between authors.
\item[\cmd{\email{address}}]
Sets the current author's email address. Only addresses of authors
marked as corresponding are printed in the correspondence block.
Literal underscores and escaped underscores are supported.
\item[\cmd{\orcid{identifier}}]
Sets the current author's ORCID identifier, without the URL prefix.
It produces an icon linking to the ORCID page. An empty argument
omits the icon. The identifier is not validated by the class.
\item[\cmd{\corresponding}]
Marks the current author as corresponding. It takes no argument.
The class suppresses the distinguishing asterisk when all authors
are corresponding. The correspondence block requires a nonempty
email address from at least one corresponding author.
\end{description}

The commands \cmd{\authorname}, \cmd{\affiliation}, \cmd{\email},
\cmd{\orcid}, and \cmd{\corresponding} require a preceding
\cmd{\author}; otherwise, the class reports an error.

\subsection{Abstract, keywords, and metadata}
\begin{description}
\item[\texttt{abstract} environment]
Collects the abstract for the next title block. Use
\cmd{\begin{abstract}} and \cmd{\end{abstract}} before
\cmd{\maketitle}. A later abstract replaces an earlier one.
\item[\cmd{\keywords{list}}]
Sets the keyword text; use semicolons or another desired separator.
\item[\cmd{\pacs{list}}]
Sets PACS classification text. It is printed after the keywords.
\item[\cmd{\arxivnumber{identifier}}]
Sets an arXiv identifier without an \texttt{arXiv:} prefix or URL.
Examples are \texttt{2507.22660}, \texttt{2603.12466v2}, and
\texttt{gr-qc/0307012}. The class checks identifier syntax, not
whether the record exists. Invalid input produces a warning and
preserves any earlier valid identifier. The link appears in the
suggested citation.
\item[\cmd{\maketitle}]
Starts a paper, first finishing the preceding pages and floats.
\item[\cmd{\appendix}]
Switches the current paper to appendix section numbering.
\end{description}

Keywords, PACS codes, and the arXiv identifier are optional.
Empty or space-only keyword and PACS arguments omit their headings.

\subsection{Publication identifiers}
\begin{description}
\item[\cmd{\ijgtppublished{year}{volume}{issue}{article}}]
Sets journal publication metadata and enables the journal citation
suggestion, publication running head, and DOI footer. For example:
\par\noindent\cmd{\ijgtppublished{2026}{2}{1}{5}}
\par\noindent This constructs
\texttt{10.53941/ijgtp.2026.100005}. The DOI consists of the year,
issue, and article number padded to at least five digits; the volume
does not enter this identifier. Use actual assigned publication data.
Arguments must contain decimal digits only. The year must be at least
2025, and the other values must be positive. Each value must fit TeX's
integer range. Leading zeros are accepted and normalized. Invalid
input produces a warning and preserves any earlier valid values.
\item[\cmd{\receiveddate{date}}]
Sets the received-date text.
\item[\cmd{\reviseddate{date}}]
Sets the revised-date text.
\item[\cmd{\accepteddate{date}}]
Sets the accepted-date text.
\item[\cmd{\publisheddate{date}}]
Sets the publication-date text.
\end{description}

Date text is printed as supplied: no date parsing or chronological validation is performed.
Without publication metadata, the running head indicates a manuscript prepared for the journal.
The copyright year is the supplied publication year, if available; otherwise, it is the current compilation year.

\subsection{Bibliography commands and formatting hooks}
\begin{description}
\item[\cmd{\bibliography{database}}]
Prints the current bibliography using the named \texttt{.bib} database,
then closes that bibliography unit. Omit the filename extension;
multiple database names may be comma-separated.
\item[\cmd{\bibliographystyle{style}}]
Selects the current paper's BibTeX style. Place it after
\cmd{\maketitle} and before \cmd{\bibliography}. It may follow citations.
\item[\cmd{\bibnumfmt{number}}]
The class redefines this \texttt{natbib} formatting hook to print a
bibliography number followed by a period. Authors normally do not call it.
\item[\cmd{\bibsection}]
The class redefines this \texttt{natbib} hook to create the unnumbered
references heading and its table-of-contents entry.
\end{description}

Use the loaded \texttt{natbib} commands, such as \cmd{\cite{key}},
\cmd{\cite{first,second}}, and \cmd{\nocite{key}}, in the usual way.
The class uses numeric citations with sorting and compression.
\cmd{\nocite{*}} includes all entries from the selected database.

\subsection{Additional table column types}
The class defines \texttt{L\{width\}}, \texttt{C\{width\}}, and
\texttt{R\{width\}}: fixed-width, vertically centered paragraph columns
with left, centered, and right horizontal alignment, respectively.
For example:
\begin{verbatim}
\begin{tabular}{L{4cm}C{2cm}R{2cm}}
Quantity & Symbol & Value \\
Decay time & $\tau$ & 1.0 \\
\end{tabular}
\end{verbatim}
The public helper \cmd{\PreserveBackslash{alignment}} applies an alignment
command while retaining the meaning of the table row terminator
\cmd{\\}. The three column types use it internally; ordinary tables
do not need to call it directly.

\end{document}
